\chapter{Deployment and Go-Live}\label{ch:deploy}

RecaLL is a local-first desktop application. Unlike a web service, it is not deployed to a
public server: each student runs the complete system on their own computer, and only the AI
requests leave the machine (and only when a cloud provider is selected). This chapter therefore
describes three forms of delivery that exist in the repository: the development launch used by
the team, the packaged desktop installer produced by the release pipeline, and an optional
container setup kept for a future hosted version.

\section{Deployment Environment}

\Cref{fig:deploydiag} in Appendix~E shows the deployment view. On the student's computer the
Tauri shell hosts the React interface in the operating system's web view and starts the FastAPI
backend as a local process on \texttt{127.0.0.1:8000}. The backend stores all data in SQLite
files and a local file store. External nodes are limited to the AI provider chosen by the student
(Gemini, OpenAI or Groq over HTTPS, or a local Ollama server on port 11434) and two optional
services: Supabase for cloud sync and Langfuse for tracing.

\begin{table}[htbp]
\centering
\caption{Deployment environments}\label{tab:deployenv}
\small
\begin{tabularx}{\textwidth}{@{}l X X@{}}
\toprule
\textbf{Environment} & \textbf{Composition} & \textbf{Purpose / status}\\
\midrule
Development & \texttt{run.bat}: uvicorn with \texttt{--reload} on 127.0.0.1:8000 and the Vite dev server (\texttt{pnpm dev}) on port 5173; \texttt{stop.bat} stops both & Daily development and the demonstrations in this report. Working.\\
Packaged desktop & Tauri~2 bundle (\texttt{productName} ``Recall'', version 0.1.5) with the backend frozen by PyInstaller as the sidecar binary \texttt{recall-backend-<target>} & End-user installation. A GitHub pre-release ``Prelaunch Version 0.1.5'' (tag \texttt{v0.1.5}, 24 August 2026) provides the installer (\Cref{sec:deploy-process}).\\
Container (optional) & \texttt{Dockerfile} (Python 3.12-slim, non-root user UID 10001, health check) and \texttt{docker-compose.yml} (PostgreSQL 16 + pgvector, backend, Nginx gateway, optional Ollama) & Prepared for a future hosted version; the desktop product uses SQLite.\\
\bottomrule
\end{tabularx}
\end{table}

\paragraph{Configuration.} Settings are read through \texttt{pydantic-settings} from environment
variables and an optional \texttt{.env} file (\texttt{app/core/config.py}). Secrets such as
\texttt{SECRET\_KEY} and \texttt{BYOK\_ENCRYPTION\_KEY} are never committed. Personal provider
keys entered in the Settings page are stored in the v1 BYOK vault encrypted with AES-256-GCM
(\Cref{sec:security}). No secret values are reproduced in this report.

\paragraph{Minimum client hardware.} Any 64-bit Windows computer able to run a modern web view
with 8\,GB of RAM is sufficient for cloud providers. Running a local model through Ollama needs
more memory and, ideally, a GPU (the project's deployment guide recommends 8\,GB RAM on CPU or a
GPU with 8\,GB VRAM).

\section{Deployment Process and CI/CD Implementation}\label{sec:deploy-process}

\subsection{Running from source (used for this report)}
\begin{enumerate}
  \item Clone the repository: \texttt{git clone https://github.com/Saimun-jd/RecallAI.git}.
  \item Create a virtual environment and install the backend:
        \texttt{python -m venv .venv} and \texttt{.venv\textbackslash Scripts\textbackslash pip install -r requirements.txt}.
  \item Install the frontend once: \texttt{cd desktop \&\& pnpm install} (Node.js 18+ and pnpm).
  \item Start everything with \texttt{run.bat}, which opens the backend and the interface.
  \item In \emph{Settings $\rightarrow$ AI \& Providers}, add a personal API key or select Ollama.
\end{enumerate}

\subsection{Building the desktop installer}
The release pipeline in \texttt{.github/workflows/release.yml} (\Cref{sec:cicd}) runs when a tag
matching \texttt{v*} is pushed, or by hand. On a \texttt{windows-latest} runner it:
(1)~installs Python~3.12 and the requirements together with PyInstaller;
(2)~freezes the backend into a single executable with \texttt{scripts/build\_sidecar.py}, which
names it with the Rust target triple and places it in \texttt{desktop/src-tauri/binaries/};
(3)~installs pnpm~9 and Node.js~20 and builds the interface;
(4)~installs the stable Rust toolchain; and
(5)~runs \texttt{tauri-apps/tauri-action}, which compiles the Tauri shell, bundles the sidecar
declared in \texttt{tauri.conf.json} (\texttt{externalBin}) and produces the Windows installer.
The workflow builds the installer but does not publish it; the \texttt{v0.1.5} pre-release, with
\texttt{Recall\_0.1.5\_x64-setup.exe}, was published on GitHub by \texttt{Saimun-jd}.

\section{Go-Live Strategy and Post-Deployment Monitoring}

\paragraph{Rollout.} Because every installation is independent, a phased rollout fits best:
\begin{enumerate}
  \item \textbf{Internal use} by the four team members on their own course material (done during
  development).
  \item \textbf{Pilot} with six CUET students (P1--P6, including one team member), who tested the
  application informally during development; this also served as the user acceptance test in
  \Cref{sec:uat}.
  \item \textbf{Public pre-release} through GitHub Releases, marked as a pre-release so that users
  know the software is still changing.
\end{enumerate}

\paragraph{User support.} The repository README and the in-app Settings page explain how to add a
provider key or use Ollama. During the pilot, support was given in person at the team's weekly
meetings in the university hall.

\paragraph{Monitoring.} There is no central server, so monitoring is local and privacy-preserving:
\begin{itemize}
  \item health endpoints \texttt{/health}, \texttt{/health/db} and \texttt{/health/ready} report
  the state of the backend and its database;
  \item every document processing job records its stage, progress and error in a job table, which
  the interface shows to the student;
  \item the token-usage log (\texttt{/api/token-usage/summary} and \texttt{/history}) lets a
  student see how many AI tokens each feature has used;
  \item optional Langfuse tracing can be enabled with the student's own keys to inspect prompts and
  latency.
\end{itemize}

\paragraph{Rollback.} Each release is a separate installer, so the student can reinstall the
previous version from GitHub Releases. The databases and uploaded files live in the user's
application-data folder (\texttt{platformdirs.user\_data\_dir("Recall")}) and are kept. Schema changes
applied at start-up are additive only: \texttt{app/models/schema\_init.py} (run by
\texttt{app/core/migrations.py}) and \texttt{init\_db()} in \texttt{app/database.py} use
\texttt{CREATE TABLE IF NOT EXISTS} and \texttt{ALTER TABLE \ldots{} ADD COLUMN} with nullable or
defaulted columns, never drop or rename anything, and have no down migrations. An older version can
therefore still open an upgraded database and ignores what it does not know; copying the data folder
before changing version is still advised.
